\documentclass[11pt,a4paper]{article}
% preamble.tex — estilo común de todos los apuntes Froth.
% Cada apunte hace % preamble.tex — estilo común de todos los apuntes Froth.
% Cada apunte hace % preamble.tex — estilo común de todos los apuntes Froth.
% Cada apunte hace \input{preamble} justo después de \documentclass.
\usepackage[margin=2.4cm]{geometry}
\usepackage{xcolor}
\definecolor{litblue}{RGB}{31,79,121}
\definecolor{litgreen}{RGB}{27,94,32}
\definecolor{litorange}{RGB}{230,81,0}
\usepackage{parskip}
\usepackage{titlesec}
\titleformat{\section}{\Large\bfseries\color{litblue}}{\thesection}{0.6em}{}
\titleformat{\subsection}{\large\bfseries\color{litblue}}{\thesubsection}{0.6em}{}
\usepackage{tcolorbox}
\usepackage{fancyhdr}
\pagestyle{fancy}
\fancyhf{}
\fancyhead[L]{\small\color{litblue}Froth — Apuntes de ML}
\fancyhead[R]{\small\thepage}
\renewcommand{\headrulewidth}{0.4pt}
\usepackage[colorlinks=true,linkcolor=litblue,urlcolor=litblue]{hyperref}

% Cabecera de cada apunte: \cabecera{numero}{titulo}
\newcommand{\cabecera}[2]{%
  \begin{tcolorbox}[colback=litblue,coltext=white,colframe=litblue,arc=2mm]
    {\Large\bfseries Apunte #1 — #2}\\[3pt]
    {\small Proyecto Froth · aprender Machine Learning construyendo}
  \end{tcolorbox}\vspace{0.8em}}

% Cajas temáticas
\newtcolorbox{concepto}[1]{colback=blue!4,colframe=litblue,title={Concepto: #1},arc=1.5mm}
\newtcolorbox{practica}{colback=green!5,colframe=litgreen,title={Pruébalo tú (teclado en mano)},arc=1.5mm}
\newtcolorbox{ojo}{colback=orange!8,colframe=litorange,title={Ojo / error típico},arc=1.5mm}
\newtcolorbox{resumen}{colback=gray!8,colframe=gray!60,title={En una frase},arc=1.5mm}

\newcommand{\codigo}[1]{\texttt{#1}}
 justo después de \documentclass.
\usepackage[margin=2.4cm]{geometry}
\usepackage{xcolor}
\definecolor{litblue}{RGB}{31,79,121}
\definecolor{litgreen}{RGB}{27,94,32}
\definecolor{litorange}{RGB}{230,81,0}
\usepackage{parskip}
\usepackage{titlesec}
\titleformat{\section}{\Large\bfseries\color{litblue}}{\thesection}{0.6em}{}
\titleformat{\subsection}{\large\bfseries\color{litblue}}{\thesubsection}{0.6em}{}
\usepackage{tcolorbox}
\usepackage{fancyhdr}
\pagestyle{fancy}
\fancyhf{}
\fancyhead[L]{\small\color{litblue}Froth — Apuntes de ML}
\fancyhead[R]{\small\thepage}
\renewcommand{\headrulewidth}{0.4pt}
\usepackage[colorlinks=true,linkcolor=litblue,urlcolor=litblue]{hyperref}

% Cabecera de cada apunte: \cabecera{numero}{titulo}
\newcommand{\cabecera}[2]{%
  \begin{tcolorbox}[colback=litblue,coltext=white,colframe=litblue,arc=2mm]
    {\Large\bfseries Apunte #1 — #2}\\[3pt]
    {\small Proyecto Froth · aprender Machine Learning construyendo}
  \end{tcolorbox}\vspace{0.8em}}

% Cajas temáticas
\newtcolorbox{concepto}[1]{colback=blue!4,colframe=litblue,title={Concepto: #1},arc=1.5mm}
\newtcolorbox{practica}{colback=green!5,colframe=litgreen,title={Pruébalo tú (teclado en mano)},arc=1.5mm}
\newtcolorbox{ojo}{colback=orange!8,colframe=litorange,title={Ojo / error típico},arc=1.5mm}
\newtcolorbox{resumen}{colback=gray!8,colframe=gray!60,title={En una frase},arc=1.5mm}

\newcommand{\codigo}[1]{\texttt{#1}}
 justo después de \documentclass.
\usepackage[margin=2.4cm]{geometry}
\usepackage{xcolor}
\definecolor{litblue}{RGB}{31,79,121}
\definecolor{litgreen}{RGB}{27,94,32}
\definecolor{litorange}{RGB}{230,81,0}
\usepackage{parskip}
\usepackage{titlesec}
\titleformat{\section}{\Large\bfseries\color{litblue}}{\thesection}{0.6em}{}
\titleformat{\subsection}{\large\bfseries\color{litblue}}{\thesubsection}{0.6em}{}
\usepackage{tcolorbox}
\usepackage{fancyhdr}
\pagestyle{fancy}
\fancyhf{}
\fancyhead[L]{\small\color{litblue}Froth — Apuntes de ML}
\fancyhead[R]{\small\thepage}
\renewcommand{\headrulewidth}{0.4pt}
\usepackage[colorlinks=true,linkcolor=litblue,urlcolor=litblue]{hyperref}

% Cabecera de cada apunte: \cabecera{numero}{titulo}
\newcommand{\cabecera}[2]{%
  \begin{tcolorbox}[colback=litblue,coltext=white,colframe=litblue,arc=2mm]
    {\Large\bfseries Apunte #1 — #2}\\[3pt]
    {\small Proyecto Froth · aprender Machine Learning construyendo}
  \end{tcolorbox}\vspace{0.8em}}

% Cajas temáticas
\newtcolorbox{concepto}[1]{colback=blue!4,colframe=litblue,title={Concepto: #1},arc=1.5mm}
\newtcolorbox{practica}{colback=green!5,colframe=litgreen,title={Pruébalo tú (teclado en mano)},arc=1.5mm}
\newtcolorbox{ojo}{colback=orange!8,colframe=litorange,title={Ojo / error típico},arc=1.5mm}
\newtcolorbox{resumen}{colback=gray!8,colframe=gray!60,title={En una frase},arc=1.5mm}

\newcommand{\codigo}[1]{\texttt{#1}}

\begin{document}
\cabecera{13}{UMAP en acción: el primer mapa de tus 126 papers}

\section{Qué hicimos}

Corrimos UMAP de verdad sobre la matriz $126 \times 768$ y salió una tabla $126 \times 2$:
una coordenada (x, y) por paper. Las guardamos en \codigo{2\_Datos/embeddings/umap\_2d.npy}
y dibujamos el primer scatter (está en \codigo{3\_Resultados/umap\_first\_look.png}).

\textbf{Lo importante del resultado:} sin decirle nada sobre minería, ya se ven
\emph{estructuras}: una nube grande a la izquierda, un brazo arriba a la derecha, una franja
diagonal abajo con un grumo apretado en la esquina. Eso son (probablemente) tus subtemas
esperando nombre. El paso siguiente (HDBSCAN) los delimitará formalmente.

\section{Las tres decisiones del código (y por qué)}

\begin{concepto}{metric="cosine" — medir con la misma vara}
UMAP necesita decidir quiénes son los ``vecinos'' de cada punto, y para eso mide distancias.
Le dijimos que las mida con \textbf{similitud coseno} — exactamente la misma vara de la
Fase 2 (Apunte 08). Si la Fase 2 compara papers por coseno y UMAP usara otra medida, el mapa
contaría una historia diferente a la del buscador. Consistencia ante todo.
\end{concepto}

\begin{concepto}{random\_state=42 — reproducibilidad}
UMAP tiene azar interno (de dónde parte el acomodo). Sin semilla fija, cada corrida daría un
mapa \emph{ligeramente} distinto — mismos grupos, otras posiciones. Fijar
\codigo{random\_state=42} hace que cualquiera que corra tu código obtenga \textbf{el mismo
mapa}. En ciencia (y en tu tesis) eso se llama \emph{reproducibilidad} y es sagrado.
(El 42 es una broma-tradición de programadores; cualquier número fijo sirve.)
\end{concepto}

\begin{concepto}{n\_components=2 — cuántas dimensiones de salida}
Le pedimos 2 dimensiones porque el objetivo es \emph{dibujar}. Se puede reducir a 5 o 10
(útil a veces para el clustering); lo exploraremos si hace falta.
\end{concepto}

\begin{ojo}
En el scatter quitamos los números de los ejes \textbf{a propósito}. Tras el Apunte 12 ya
sabes por qué: en UMAP los ejes no significan nada. Si algún día ves un mapa UMAP con ejes
numerados e interpretados (``a mayor X, mayor tal cosa''), desconfía del autor.
\end{ojo}

\section{Cómo leer tu primer mapa}

\begin{itemize}
  \item \textbf{Nubes densas} = grupos de papers que hablan de lo mismo (futuros clusters).
  \item \textbf{Puntos sueltos / puentes} = papers que mezclan temas o que no encajan
        (candidatos a ``ruido'' en HDBSCAN).
  \item \textbf{Forma general} (brazo, franja, grumo): sugiere subtemas con distinta cohesión
        --- un grumo apretado es un tema muy específico; una nube amplia, un tema general.
\end{itemize}

\begin{resumen}
UMAP convirtió 768 dimensiones en un mapa 2D reproducible (semilla fija) midiendo vecinos
con coseno (la misma vara de siempre), y las nubes aparecieron solas: son los subtemas que
HDBSCAN delimitará en el siguiente paso.
\end{resumen}

\end{document}
